\subsection{Nondegenerate cells and normal multiplicities}
\label{sec:flag-nondegenerate-normal-count}

We work with actual admitted nonzero-Gram returns, including
nongeneric ones. Write \(E/F\) for the CM host, \(f=[F:\mathbf Q]\),
\(\rho=\operatorname{rd}(E)\), and \(\mathfrak d\) for the extended
different, of absolute norm \(\rho^{2f}\). The level \(g\) is the
product of the two split rational primes, prime to \(\mathfrak d\).
Use the pairing \(B\) linear in its first variable. A return is an
isometry \(\gamma_t:\Lambda'\to\Lambda\) of the actual row lattices.
Write its target and transported source frames as
\((e_0,\dots,e_r)\), \((f_0,\dots,f_r)\), with last vectors \(z,z'\)
of norms \(-T,-T'\), isotropic spans \(X,Y_t=\gamma_t X'\), and
signature \((k+r,r+1)\). Put
\[
 \begin{gathered}
 w_a=r-a,\qquad w=r(r+1),\qquad w_l=\sum_{a<l}w_a,\\
 Z_{ab}=B(f_b,e_a),\qquad D_l=\det Z_{[0,l),[0,l)},\\
 D_0=1,\qquad G=D_{r+1}\bar D_{r+1}-TT'D_r\bar D_r.
 \end{gathered}
\]
The actual weighted lattice conditions give
\(e_a,f_a\in\mathfrak d^{w_a}\Lambda\).
Above \(g\), admission says that the full second frame is the
first frame times an integral upper unitriangular matrix modulo
\(g\Lambda_{\mathfrak p}\).

\begin{lemma}[Occupied nonzero-Gram points]
If \(G\ne0\), then \(G\) is totally positive. For a relevant term,
\(Z\operatorname{diag}((-1)^{r-b})\) lies in the unique original
block-antidiagonal point cell under lower and upper unitriangular
multiplication over \(E\).
\end{lemma}
\begin{proof}
The Gram determinant of the concatenated frames is
\((-1)^{r+1}G\): in the conjugate-linear-first convention its
diagonal blocks are \(-TE_{rr},-T'E_{rr}\), and its upper cross
block is \(Z\). Swapping block columns and taking the rank-one
Schur determinant proves the identity for invertible \(Z\),
hence identically. The joint span \(J_t\) is nondegenerate of
dimension \(2r+2\) and has negative index \(r+1\), since it
contains the first isotropic \(r\)-plane and its perpendicular
negative vector. Its complement \(H_t\) is positive of rank
\(k-1\). The determinant sign gives total positivity of \(G\),
so \(D_{r+1}\ne0\).

If \(A,A'\) are strictly upper triangular with
\(A^*Z+ZA'=0\), acting by these on the two frame spans and by
zero on \(H_t\) gives a skew-Hermitian endomorphism. It belongs
to both flag Lie algebras: it lowers the flags; in \(X^\perp\)
the second-frame component lies in the one-dimensional kernel
of the first \(r\) rows of \(Z\). That line is invariant under
\(A'\) by skew-Hermitianity, so nilpotence makes its action zero.
The reverse argument applies to \(Y_t\). Relevance equates the
traces of the immediate-superdiagonal sums \(s(A),s(A')\).
Applying relevance also to \(dA,\bar d A'\) for every \(d\in E\)
gives \(s(A)=\overline{s(A')}\) by trace nondegeneracy.
The signed monomial reduction argument of
Lemma~\ref{lem:rational-flag-signed-block}, now in size \(r+1\),
therefore applies. The permutations reverse consecutive blocks,
with alternating signs removed by the displayed column signs.
Uniqueness of the monomial representative in a
lower-unitriangular/upper-unitriangular double cell follows by
successive pivot elimination (the first nonzero position and
value in the first column are invariant, and one continues on
the remaining ordered rows and columns). These operations adapt
the two flags, preserving their graded vectors and last-vector
classes. The leading minors are unchanged. Changing the original
return representative by either flag unipotent likewise
only adapts the frames.
\end{proof}

\begin{lemma}[Full level content of the two normals]
\label{lem:nonzero-gram-normal-content}
For every \(l\), \(D_l\in\mathfrak d^{2w_l}\). For \(1\le l\le r\),
\(D_l\in(g)^l\), and \(D_{r+1}\in(g)^r\). Moreover
\[
 G\in(g)^{2(r+1)}\mathfrak d^{2w}.
 \tag{NG1}
\]
Define the source's adjugate normal vectors by
\[
 \omega=\sum_b(\operatorname{adj}Z)_{br}f_b,\qquad
 y=T\omega+D_{r+1}z,\qquad
 y'=\gamma_t^{-1}y(F',F),
\]
where the reverse-pair construction uses \(T'\).
They are perpendicular to their respective isotropic spans
and last vectors and have norms \(TG,T'G\).
Their quotient classes belong to
\(\mathfrak c Q_\Lambda,\mathfrak c Q_{\Lambda'}\), where
\[
 Q_\Lambda=(X^\perp\cap\Lambda)/(X\cap\Lambda),\qquad
 Q_{\Lambda'}=((X')^\perp\cap\Lambda')/(X'\cap\Lambda'),
 \qquad \mathfrak c=(g)^{r+1}\mathfrak d^w.                 \tag{NG2}
\]
\end{lemma}
\begin{proof}
The weighted lattice conditions give cross-entry content
\(\mathfrak d^{w_a+w_b}\); determinant expansion proves the
different assertions. Above \(g\), the first \(r\) cross rows
vanish modulo \(g\). Subtract the first frame times the
admission matrix from the second frame in the joint Gram
determinant: the last \(r+1\) rows and columns of the new
matrix each have a factor \(g\). This proves (NG1) by
coprimality of \(g,\mathfrak d\), and the cross-row condition
gives the other minor divisibilities.

The adjugate identity gives \(B(\omega,e_a)=0\) for \(a<r\)
and \(B(\omega,z)=D_{r+1}\). Its norm is
\(-T'D_r\bar D_r\), by the second frame's internal Gram
matrix. These give the asserted properties of \(y\), and
reversal gives those of \(y'\). The cofactor omitting row
\(r\) and column \(b\) lies in \(\mathfrak d^{w-w_b}\);
multiplying by \(f_b\) shows \(y\in\mathfrak d^w\Lambda\).
At a prime above \(g\), write the cofactors as \(a_b\in(g)^r\).
Admission gives \(f_b\in X+g\Lambda_{\mathfrak p}\) for \(b<r\)
and \(z'\in z+X+g\Lambda_{\mathfrak p}\). Thus
\[
 \omega\equiv a_rz\pmod{X+g^{r+1}\Lambda_{\mathfrak p}},
 \qquad D_{r+1}=\sum_b Z_{rb}a_b\equiv-Ta_r\pmod{g^{r+1}} .
\]
Therefore \(y\in X+g^{r+1}\Lambda_{\mathfrak p}\).
Perpendicularity to \(X\) makes the lattice vector in this
expression also perpendicular to \(X\). Reversing the
admission using the inverse unitriangular matrix, then
transporting by \(\gamma_t^{-1}\), gives the same source
condition. Local intersection proves (NG2).
\end{proof}

In the actual point normalization put
\[
 N_G=\prod_\tau\rho^{-2w}\tau(G),\qquad
 N_D=\prod_{l\ {\rm block\ end}}\prod_\tau
              \rho^{-4w_l}|\tau(D_l)|^2,\qquad
 G_0=g^{-2f(r+1)}N_G.
\]
Here \(\tau\) ranges over real places of \(F\) with a chosen
complex extension. Each endpoint factor is the norm of the
integral ideal \((D_l)\mathfrak d^{-2w_l}\).
The ideal \((G)(g)^{-2(r+1)}\mathfrak d^{-2w}\) is integral
and extended from \(F\); its \(F\)-norm is \(G_0\), a positive
integer. In particular these are the actual occupied-point
StepPointsG and StepNormOneG inequalities, strengthened to
\(N_G\ge g^{2f(r+1)}\). The original height factor
\(g^{-(2r+1)f}N_GN_D\) is at most \(g^fN_GN_D\).

\begin{theorem}[Level-saving rational normal count]
\label{thm:nonzero-gram-normal-count}
Fix an actual row pair, a corank-one class of nonzero-Gram
returns and an occupied point. The number of rational
unipotent double cosets in that class and point is at most
\[
             |\mu(E)|^2(TT')^fG_0^2 .                     \tag{NG3}
\]
This proves the occupied normal-count input for the generic
scalar classes and the nongeneric corank-one classes, with
the source's larger bound
\(\rho^{12wf}(TT')^{6f}N_G^6\) replaced by (NG3).
It does not count integral double cosets.
\end{theorem}
\begin{proof}
Put \(P=(X^\perp\cap z^\perp)/X\), and \(P'\) for the fixed
source flag. Both are positive of rank \(k\). The two
isometric injections of \(H_t\) into \(P,P'\), the latter
via \(\gamma_t^{-1}\), define a graph \(\mathcal H_t\)
of dimension \(k-1\) in \(P\oplus P'\).
The source's ThPairs conditions
\(h_1-h_2\in X\) and
\(\gamma_2\gamma_1^{-1}h_1-h_2\in Y_{t_2}\) say precisely
that the two graph vectors agree. Both projections of
each graph are injective. Hence
\[
 \dim\operatorname{Th}=\dim(\mathcal H_{t_1}\cap\mathcal H_{t_2}),
 \quad s_0=1\ \Longleftrightarrow\
                  \mathcal H_{t_1}=\mathcal H_{t_2}.
\]
This proves the corank-one equivalence relation, including
\(k=1\). Within it the generic scalar refinement is also
an equivalence: the spaces Ylift and Ylift' in the source
project onto the perpendicular lines to the two graph
projections, and SgRel records the two last-vector
coefficients of the difference of lifts in the direct
sum of the two frame spans. They are unique and independent
of lifts modulo the corresponding isotropic spaces.
On identifying the two perpendicular lines by the fixed
quotient coordinates \(P,P'\), IsScalar says exactly that
these coefficient functions coincide. Restriction to such
a class can only reduce the count.

Fix a base term \(t_0\) in this class and point. The two
normal classes lie on fixed perpendicular lines and have
norms \(TG,T'G\) by the preceding lemma. Write
\([y_t]=a_t[y_{t_0}], [y'_t]=a'_t[y'_{t_0}]\).
These are the source coefficient ratios obtained by
pairing with the base normal and dividing by its norm.
Each has relative norm one. The fractional ideal
\[
 I=\{a\in E:a[y_{t_0}]\in\mathfrak c Q_\Lambda\}
\]
contains all \(a_t\), and is nonzero by (NG2); the class
of the base vector is nonzero. The integral pairing on
the quotient gives \(TG I\bar I\subseteq
\mathfrak c\bar{\mathfrak c}\). Taking absolute ideal
norms gives \(N_E I\ge1/(T^f G_0)\).
Similarly \(N_E I'\ge1/((T')^fG_0)\) for the source.

To verify the rational-normal separation, suppose two
terms have identical pairs of ratios. Adapt both pairs
of frames to the same cross matrix (remove the fixed
column signs from the common signed point).
The map on their joint spans sending the adapted
frames to each other is an isometry, since their full
Gram matrices agree. Extend it by the isometry of
the two complements given by the common target graph
projection. This extension \(\sigma\) is rational and
also respects the source graph projection.

Under unitriangular adaptations \(A,B\) the cross
matrix becomes \(A^*ZB\); the adjugate vector in the
ambient space is unchanged because
\((A^*)^{-1}e_r^{\rm coord}=e_r^{\rm coord}\).
The determinant is unchanged. Hence the normal changes
only by an element of the isotropic span; the reverse
normal does likewise. The extension \(\sigma\) thus
fixes the first quotient normal class, the graph
projection and the last-vector class. These span
\(X^\perp/X\). It fixes each graded flag vector, so
\(\sigma\in N(E)\), the original flag unipotent.
On the source side
\(s'=\gamma_1^{-1}\sigma^{-1}\gamma_2\) fixes every
graded vector and acts trivially on the isotropic
quotient by the common graph and equality of source
normals. Thus \(s'\in N'(E)\), and
\(\gamma_2=\sigma\gamma_1s'\).
Choosing one representative of each occurring rational
double coset gives an injection into pairs of ratios.

For any fractional ideal \(I\) and nonzero \(m\), an
element \(a\in I\) satisfying \(a\bar a=m\) determines
the integral ideal \(A=(a)I^{-1}\), of norm
\(L=\sqrt{|N_E(m)|}/N_E I\). Its product with its
conjugate is fixed. At distinct conjugate primes with
fixed total exponent \(e\), at most \(e+1\le q^e\)
choices occur, with \(q\) their common absolute norm;
at fixed primes there is at most one choice.
Thus the number of ideals is at most \(L\). In one
ideal the elements differ by norm-one integral units,
which are roots of unity by Kronecker. This proves the
stronger norm-equation count \(|\mu(E)|L\) whenever the
solution set is nonempty. Applying it to \(I,I'\) with
\(m=1\) gives (NG3). Finally \(|\mu(E)|\le8f^2\) by
\(\varphi(n)^2\ge n/2\) and \(\varphi(|\mu(E)|)\le2f\).
For \(T,T'\ge3\) we have \(64f^4\le(TT')^{4f}\);
\(G_0\le N_G\), \(N_G\ge1\), prove the original bound.
\end{proof}

This supplies the occupied-point and rational-normal
arithmetic steps of the original nondegenerate-return
calculation. The integral amplitude sum within each
rational double coset, local generic amplitudes, and
the restricted arithmetic theta comparison remain separate.

\paragraph{Arbitrarily small exponent in the norm count.}
For every fixed \(\epsilon>0\), the norm-equation proof also gives
\[
 \#\{a\in I:a\bar a=m\}
 \le |\mu(E)|C_\epsilon^f L^\epsilon,\qquad
 C_\epsilon=\prod_{p\le 2^{1/\epsilon}}
                         (1-p^{-\epsilon})^{-2}.
 \tag{NG4}
\]
In fact \(e+1\le 2^e\le q^{\epsilon e}\) at all conjugate
prime pairs of absolute norm \(q\ge 2^{1/\epsilon}\).
At the other pairs,
\((e+1)q^{-\epsilon e}\le
\sum_{j\ge0}(j+1)q^{-\epsilon j}=(1-q^{-\epsilon})^{-2}\).
There are at most \(f\) such distinct conjugate pairs
above each rational prime \(p\), and \(q\ge p\).
Multiplication proves (NG4), where the finite constant
depends only on the prescribed epsilon, not on the field.
Thus (NG3) can further be replaced by
\[
 |\mu(E)|^2 C_\epsilon^{2f}(TT')^{\epsilon f}G_0^{2\epsilon}.
\]
This bound is on the same occupied normal classes; no
power of the integral multiplicity is being removed.
